\section*{Chapitre \ref{ch:g9:periodic-table-first-look} --- Le tableau périodique~: un premier regard}

\begin{solution}{exo:g9:periodic-table-first-look:1}
Par \omterm{def:g9:inside-the-atom:atomic-number}{numéro atomique} $Z$ croissant.
\end{solution}

\begin{solution}{exo:g9:periodic-table-first-look:2}
Le carbone~: \omterm{def:g9:periodic-table-first-look:period}{période} 2, \omterm{def:g9:periodic-table-first-look:period}{groupe} 14. Le magnésium~: \omterm{def:g9:periodic-table-first-look:period}{période} 3, \omterm{def:g9:periodic-table-first-look:period}{groupe} 2.
L'argon~: \omterm{def:g9:periodic-table-first-look:period}{période} 3, \omterm{def:g9:periodic-table-first-look:period}{groupe} 18.
\end{solution}

\begin{solution}{exo:g9:periodic-table-first-look:3}
Métaux~: le fer, le cuivre, l'aluminium. Non-métaux~: le soufre,
l'oxygène, le chlore.
\end{solution}

\begin{solution}{exo:g9:periodic-table-first-look:4}
Le potassium (\ce{K}) sous le sodium~; le brome (\ce{Br}) sous le
chlore.
\end{solution}

\begin{solution}{exo:g9:periodic-table-first-look:5}
Le calcium (\ce{Ca}, $Z = 20$) et l'azote (\ce{N}, $Z = 7$).
\end{solution}

\begin{solution}{exo:g9:periodic-table-first-look:6}
Il est dans le même groupe que le sodium, le \omterm{def:g9:periodic-table-first-look:period}{groupe} 1, et les éléments
d'un même groupe se comportent de façon semblable.
\end{solution}

\begin{solution}{exo:g9:periodic-table-first-look:7}
Un \omterm{def:g9:periodic-table-first-look:metal}{métal}~: il est brillant, conduit l'électricité et forme un \omterm{def:g9:ions:ion}{ion}
positif.
\end{solution}

\begin{solution}{exo:g9:periodic-table-first-look:8}
Pour garder dans une même colonne les éléments aux propriétés
semblables, même s'il fallait pour cela laisser une case vide. Les cases
vides étaient des prédictions~: les éléments manquants furent trouvés
plus tard, avec les propriétés qu'il avait décrites.
\end{solution}

\begin{solution}{exo:g9:periodic-table-first-look:9}
Le carbone~: $\ce{C} = 12$~; l'oxygène~: $\ce{O} = 16$. Les cases
«~? $= 68$~» (à côté de Al $= 27.4$) et «~? $= 70$~» (à côté de
Si $= 28$).
\end{solution}

\begin{solution}{exo:g9:periodic-table-first-look:10}
\omterm{def:g9:periodic-table-first-look:period}{Période} 2~: 8 éléments (du lithium au néon). \omterm{def:g9:periodic-table-first-look:period}{Période} 4~: 18 éléments
(du potassium au krypton).
\end{solution}

\begin{solution}{exo:g9:periodic-table-first-look:11}
Le \omterm{def:g9:periodic-table-first-look:period}{groupe} 18. Là où rien ne doit réagir~: l'argon des ampoules et des
gaz de soudage, l'hélium des ballons (il ne \omterm{def:g4:what-a-fire-needs:burning}{brûle} pas).
\end{solution}

\begin{solution}{exo:g9:periodic-table-first-look:12}
$(27.0 + 114.8)/2 = 70.9$, proche de 69.7.
\end{solution}

\begin{solution}{pb:g9:periodic-table-first-look:1}
\textbf{1.} $Z = 32$~; \omterm{def:g9:periodic-table-first-look:period}{période} 4, \omterm{def:g9:periodic-table-first-look:period}{groupe} 14.

\textbf{2.} Au-dessus~: le silicium. Au-dessous~: l'étain. À gauche~: le
gallium. À droite~: l'arsenic.

\textbf{3.} Un métalloïde.

\textbf{4.} Il est dans le même groupe que le silicium, et les éléments
d'un même groupe se comportent de façon semblable.

\textbf{5.} $(28.1 + 118.7)/2 = 73.4$.

\textbf{6.} $(69.7 + 74.9)/2 = 72.3$.

\textbf{7.} La moyenne gauche--droite (72.3)~: le long d'une \omterm{def:g9:periodic-table-first-look:period}{période}, la
masse atomique augmente régulièrement, d'un élément à l'autre, tandis
qu'en descendant un groupe elle fait de grands sauts.

\textbf{8.} $72.6 - 70 = 2.6$.

\textbf{9.} $(28.1 + 118.7 + 69.7 + 74.9)/4 = 291.4/4 = 72.85$, soit
72.9.

\textbf{10.} La moyenne, 72.9, est proche à la fois du 72 de Mendeleïev
et du 72.3 de Winkler (et de la valeur actuelle, 72.6).

\textbf{11.} Un élément prédit avant que quiconque l'ait vu fut trouvé
avec les propriétés annoncées~: le tableau n'était pas seulement une
liste, il permettait de faire des prédictions.

\textbf{12.} Environ 72.9.
\end{solution}
